\documentclass[9pt,a4paper]{extarticle}
\usepackage[margin=1.6cm]{geometry}
\usepackage[T1]{fontenc}
\usepackage[utf8]{inputenc}
\usepackage[colorlinks=true,urlcolor=blue]{hyperref}
\usepackage{enumitem,titlesec}
\setlength{\parindent}{0pt}
\setlist[enumerate]{leftmargin=*,itemsep=6pt,topsep=3pt,parsep=0pt}
\titlespacing*{\section}{0pt}{9pt}{4pt}
\titleformat{\section}{\normalsize\bfseries}{}{0pt}{}
\begin{document}
\begin{center}{\Large\bfseries Tabular and Time-Series Foundation Models}\\[2pt]{\small Selected bibliography --- Valentin Kilian --- October 2026}\end{center}


\section{Foundations}
\begin{enumerate}
\item \textbf{\href{https://arxiv.org/abs/2112.10510}{Transformers Can Do Bayesian Inference}}.\\
Samuel M\"uller, Noah Hollmann, Sebastian Pineda Arango, Josif Grabocka, Frank Hutter (2022). ICLR.\\
 Introduces prior-data fitted networks (PFNs).
\item \textbf{\href{https://arxiv.org/abs/2207.01848}{TabPFN: A Transformer That Solves Small Tabular Classification Problems in a Second}}.\\
Noah Hollmann, Samuel M\"uller, Katharina Eggensperger, Frank Hutter (2023). ICLR.\\
 Applies PFNs to tabular classification, enabling fast in-context predictions without dataset-specific optimisation.
\item \textbf{\href{https://doi.org/10.1038/s41586-024-08328-6}{Accurate predictions on small data with a tabular foundation model}}.\\
Noah Hollmann, Samuel M\"uller, Lennart Purucker, Arjun Krishnakumar, Max K\"orfer, Shi Bin Hoo, Robin Tibor Schirrmeister, Frank Hutter (2025). Nature.\\
 Presents TabPFNv2, adding regression and improving prediction quality through richer  priors and inference methods.
\item \textbf{\href{https://arxiv.org/abs/2502.05564}{TabICL: A Tabular Foundation Model for In-Context Learning on Large Data}}.\\
Jingang Qu, David Holzm\"uller, Gaël Varoquaux, Marine Le Morvan (2025). ICML.\\
 Introduces a two-stage architecture that makes tabular in-context learning practical on much larger datasets.
\item \textbf{\href{https://arxiv.org/abs/2602.11139}{TabICLv2: A better, faster, scalable, and open tabular foundation model}}.\\
Jingang Qu, David Holzm\"uller, Gaël Varoquaux, Marine Le Morvan (2026). arXiv.\\
 Combines an improved synthetic training generator with efficient attention mechanisms for regression and classification.
\end{enumerate}
\section{Other tabular fondations models}
 \begin{enumerate}[resume]
\item \textbf{\href{https://arxiv.org/abs/2511.08667}{TabPFN-2.5: Advancing the State of the Art in Tabular Foundation Models}}.\\
L\'eo Grinsztajn, et al. (2025). arXiv.\\
 Scales the TabPFN family to larger tables and improves zero-shot performance relative to earlier versions.
\item \textbf{\href{https://arxiv.org/abs/2605.13986}{TabPFN-3: Technical Report}}.\\
L\'eo Grinsztajn, et al. (2026). arXiv.\\
 Further increases scale to datasets with up to one million rows while reducing training and inference costs.
\item \textbf{\href{https://arxiv.org/abs/2609.17895}{TabPFN-3.5: Technical Report}}.\\
Benjamin J\"ager et al. (2026). arXiv.\\
 Extends TabPFN to more realistic settings, including grouped and temporal splits and heterogeneous feature types.
\item \textbf{\href{https://arxiv.org/abs/2609.37959}{TabFM: A Zero-Shot Foundation Model for Tabular Data}}.\\
Weihao Kong, Erez Louidor Ilan, Shuxin Nie, Taman Narayan, Rajat Sen, Yichen Zhou, Deqing Fu, Samet Oymak, Abhimanyu Das (2026). arXiv / Google Research.\\
 Presents Google’s zero-shot tabular ICL model and its TabFM+ inference-time enhancement.
\item \textbf{\href{https://arxiv.org/abs/2609.37989}{TabFM-Auto: Self-Evolving Pipelines for Tabular Foundation Models}}.\\
Deqing Fu, Huangyuan Su, Rajat Sen, Taman Narayan, Sujay Sanghavi, Abhimanyu Das, Weihao Kong (2026). arXiv / Google Research.\\
 Uses an LLM agent to evolve data-cleaning and feature-engineering pipelines around a frozen TabFM model.
\end{enumerate}
\section{Time-series foundation models}
\begin{enumerate}[resume]
\item \textbf{\href{https://arxiv.org/abs/2310.10688}{A decoder-only foundation model for time-series forecasting}}.\\
Abhimanyu Das, Weihao Kong, Rajat Sen, Yichen Zhou (2024). ICML.\\
 Introduces TimesFM, a pretrained patch-based decoder-only transformer for zero-shot forecasting.
\item \textbf{\href{https://github.com/google-research/timesfm}{TimesFM 2.0, 2.5 and 3.0 (software releases)}}.\\
Google Research (2025). Model documentation / repository.\\
 Records the subsequent TimesFM model releases; this is a repository reference rather than a separate research paper.
\item \textbf{\href{https://arxiv.org/abs/2403.07815}{Chronos: Learning the Language of Time Series}}.\\
Abdul Fatir Ansari et al. (2024). TMLR.\\
 Tokenises continuous time-series values so pretrained language-model architectures can produce probabilistic forecasts.
%\item \textbf{\href{https://github.com/amazon-science/chronos-forecasting}{Chronos-Bolt (model release)}}.\\
Amazon Science (2024). Model repository.\\
 Moves towards faster direct multi-step forecasting instead of autoregressive token generation.
\item \textbf{\href{https://arxiv.org/abs/2501.02945}{From Tables to Time: Extending TabPFN-v2 to Time Series Forecasting}}.\\
Shi Bin Hoo, Samuel M\"uller, David Salinas, Frank Hutter (2025). arXiv (revised 2026).\\
 Shows that tabular models can forecast time series using simple temporal feature engineering without time-series pretraining.
\item \textbf{\href{https://arxiv.org/abs/2505.23719}{TiRex: Zero-Shot Forecasting Across Long and Short Horizons with Enhanced In-Context Learning}}.\\
Andreas Auer, Patrick Podest, Daniel Klotz, Sebastian B\"ock, G\"unter Klambauer, Sepp Hochreiter (2025). NeurIPS.\\
 Explores xLSTM-based in-context forecasting as an alternative to transformer time-series foundation models.
\item \textbf{\href{https://arxiv.org/abs/2510.15821}{Chronos-2: From Univariate to Universal Forecasting}}.\\
Abdul Fatir Ansari et al. (2025). arXiv.\\
 Extends Chronos to zero-shot multivariate and covariate-informed forecasting using group attention.
\end{enumerate}
\section{Benchmarks}
\begin{enumerate}[resume]
\item \textbf{\href{https://arxiv.org/abs/2506.16791}{TabArena: A Living Benchmark for Machine Learning on Tabular Data}}.\\
Nick Erickson, Lennart Purucker, Andrej Tschalzev, David Holzm\"uller, Prateek Mutalik Desai, David Salinas, Frank Hutter (2025). NeurIPS.\\
 Introduces a continuously maintained benchmark for comparing tabular foundation models with strong traditional baselines.
\item \textbf{\href{https://arxiv.org/abs/2606.30410}{Beyond IID: How General Are Tabular Foundation Models, Really?}}.\\
Lennart Purucker, Andrej Tschalzev, Nick Erickson, Gioia Blayer, David Holzm\"uller, Alan Arazi, Alexander Pfefferle, Mustafa Tajjar, Gaël Varoquaux, Frank Hutter (2026). arXiv (BeyondArena).\\
 Tests tabular models beyond IID assumptions on temporal, grouped, high-dimensional and heterogeneous datasets.
\end{enumerate}
\end{document}